\lecture{18}{Tue 25 Nov 2025 12:33}{Gravitational waves}

Recall we are working in a metric $g_{\mu \nu } =\eta _{\mu \nu } +h_{\mu \nu } $ for $h$ small, meaning we take $h^2=0$. Also typo above: $h_{00} =-2\Phi $.

Recall the gauge theory of E\&M:
\begin{mdframed}
	\[
		A_\mu \to A_\mu +\partial _\mu \lambda (x)
	.\]
\end{mdframed}
The gauge transformation for gravity is
\[
	h_{\mu \nu } \to h_{\mu \nu } -\partial _\mu \varepsilon _\nu -\partial _\nu \varepsilon _\mu 
.\]
Componentwise,
\[
	h_{00} \to h_{00} -2 \partial _0\varepsilon _0
\]
\[
	\implies \Phi \to \Phi +\partial _0\varepsilon _0
.\]
So we obtain
\[
	w_i \to w_i-\partial _0\varepsilon^i+\partial _i\varepsilon ^0
.\]
\[
	\Psi \to \Psi +\frac{1}{3} \partial _i \varepsilon ^i
.\]
\[
	s_{ij} \to \partial _i\varepsilon _j+\frac{1}{2} \partial _t\varepsilon ^t\delta _{ij} 
.\]
Since $e^\mu $ has 4 components, there are 4 diffeomorphisms. Each gauge transformation removes a degree of freedom, so we lose 4 more degrees of freedom, leaving us with 2. 

We can choose $\varepsilon ^\mu (x)^2 $ to make $G_{\mu \nu } =0$. Gravitational waves are fluctuations in a vaccuum. We will impose $\partial _is^{ij} =0$ and $\partial _iw^i=0$. This reduces our Einstein tensor
\[
	G_{00} =2\nabla ^2 \Psi =0,\qquad G_{0i} =-\frac{1}{2} \nabla ^2 s_{i} +   \partial _0 \partial _i \Psi =0
.\]
From the first equation we obtain $\Psi =0$, and from the second we obtain $\omega =0$. Our equation is now much simpler:
\[
	G^i\!_i=2\nabla ^2\Phi 
\]
\[
	G^i\!_j=(\delta ^i_j\nabla ^2-\partial _i \partial _j)\Phi +(\partial _0^2-\nabla ^2)s_{ij} =0
.\]
By definition $s_{ij} $ is traceless, so
\[
	G^i\!_i = 2\nabla ^2\Phi  \implies \Phi =\frac{1}{\nabla ^2} 0 = 0
.\]
Therefore $\left( \partial _0^2 - \nabla ^2 \right) s_{ij} =0$. Recall that since $\Phi =\Psi =\omega ^i=0$,
\[
	h_{\mu \nu } =\begin{pmatrix}
	0 & 0 & 0 & 0 \\
	0 & s^{11}  & s^{12}  & s^{13}  \\
	0 & s^{21}  & s^{22}  & s^{23}  \\
	0 & s^{31}  & s^{32}  & s^{33} 
	\end{pmatrix}
.\]
This is easier to use if we pass to Fourier space.
\[
	\widetilde{s} _{ij} =\int e^{ik_\mu x^\mu } s_{ij} (x)\,\mathrm{d} ^4x
.\]
We can Fourier transform both sides.
\[
	\int e^{ik_\mu x^\mu } \left( \partial _0^2 - \nabla ^2 \right) s_{ij}  \,\mathrm{d} ^4x
.\]
I will trust the professor that we obtain
\[
	k_0^2-k^2=0
.\]
Then
\[
	k_\mu k^\mu =0
.\]
This is like the 4-momentum, and in momentum,
\[
	p_0^2 - (\vec{p}) ^2=m^2_{\text{rest} } 
.\]
Thus the graviton is massless. Since the only degree of freedom left over is $s_{ij} $ with spin 2, gravitational waves are massless of spin 2. There is a fascinating implication we will not cover which is that if we assume there exists a massless spin 2 particle, then we obtain Einstein's entire theory of gravity.





\[
	g_{\mu \nu } =\begin{pmatrix}
	1 &  \\
	 & R
	\end{pmatrix}
.\]

